\scanpage{34}
\unsure{Cela|Ceci|On} prouve (ii). Pour (i), on note qu'on sait
déjà que si $Z$ est alg. équiv. à 0, alors
l'hom correspondant $u:A_Y(k)\longrightarrow A^i_{\mathrm{alb}}(X)$ est nul
(th. 1.6.), donc si $Z$ est $\tau$-équivalent à 0, alors
l'hom $u$ est tel que $mu=0$ pour quelque
$m>0$, donc $u=0$ car $A_Y(k)$ est $m$-divisible.

\uline{Corollaire} \unsure{1.7.|2.2.|2.7.} Pour qu'un cycle algébrique
de codim $i$ sur $X$ soit alg. équiv. à 0, il
f. et s. qu'il soit Albanese-équivalent à
un diviseur de la forme $Z(a_0)$, où
\deleted{$A$ est une} $Z$ est un cycle de \unsure{diodim|dicodim|dicdim|codim} $i$
sur un \deleted{VA $A$,} produit $\underline{A}\times X$, $A$ une
variété abélienne, tel que $\deleted{\unsure{A(k)|\mathrm{Alb}}}a\leadsto Z(a):$
$A\longrightarrow\deleted{\unsure{Z^i(}}A^i_{\mathrm{alb}}(X)$ soit un hom de
groupes.

\uline{Corollaire} 2.3. Pour tt entier $n>0$, le ss-groupe
$P^i(X)$ \unsure{des|de} \deleted{classes d'équivalence, à équiv. d'Alb.} \inserted{$A^i_{\mathrm{alb}}(X)$, formé des classes de cycles}
alg. équiv. à 0, est $n$ \uline{divisible}.

\uline{Corollaire} 2.4. Soient $Z,Z'\in A^*_{\mathrm{alb}}(X)$, dont l'un est
alg. équiv. à 0 et l'autre $\tau$-équiv. à 0.
Alors $ZZ'=0$.
% archive sheet number: 33
